\subsection{线性形式的指数}

设 V 是有限维向量空间，\(\mathbf{S}(V)\) 是其对称代数。\(\mathbf{S}(V)\) 的余代数结构在 \(\mathbf{S}(V)^{*}\) 上定义了交换且结合的代数结构（《代数》，第 III 章，§11，第 579 至 582 页）。向量空间 \(\mathbf{S}(V)^{*}\) 可典则地识别为 \(\prod_{m\geq0}\mathbf{S}^{m}(V)^{*}\)，而 \(\mathbf{S}^{m}(V)^{*}\) 可典则地识别为 V 上的对称 m-线性形式空间。将 \(V^{*}=\mathbf{S}^{1}(V)^{*}\) 典则嵌入 \(\mathbf{S}(V)^{*}\) 中，定义了从代数 \(\mathbf{S}(V^{*})\) 到代数 \(\mathbf{S}(V)^{*}\) 的单射同态，其像为 \(\mathbf{S}(V)^{*gr}=\sum_{m\geq0}\mathbf{S}^{m}(V)^{*}\)（《代数》，第 III 章，§11，第 5 节，命题 8）。我们通过此同态识别代数 \(\mathbf{S}(V^{*})\) 与 \(\mathbf{S}(V)^{*gr}\)；还将 \(\mathbf{S}(V^{*})\) 识别为 V 上多项式函数的代数（第 VII 章，附录 I，第 1 节）。

满足 \((u_m) \in \prod_{m \geq 0} \mathbf{S}^m(\mathrm{V})^*\) 的元素 \(u_0 = 0\) 构成 \(\mathbf{S}(\mathrm{V})^*\) 的一个理想 J；我们在 \(\mathbf{S}(\mathrm{V})^*\) 上赋予 J-进拓扑（《交换代数》，第 III 章，§2，第 5 节），在此拓扑下 \(\mathbf{S}(\mathrm{V})^*\) 是完备的，而 \(\mathbf{S}(\mathrm{V}^*)\) 在 \(\mathbf{S}(\mathrm{V})^*\) 中稠密。若 \((e_i^*)_{1 \leq i \leq n}\) 是 \(\mathrm{V}^*\) 的一个基，且 \(\mathrm{T}_1, \ldots, \mathrm{T}_n\) 是不定元，则从 \(k[\mathrm{T}_1, \ldots, \mathrm{T}_n]\) 到 \(\mathbf{S}(\mathrm{V}^*)\) 的、将 \(\mathrm{T}_i\) 映到 \(e_i^*(1 \leq i \leq n)\) 的同态是代数同构，并延拓为从代数 \(k[[\mathrm{T}_1, \ldots, \mathrm{T}_n]]\) 到代数 \(\mathbf{S}(\mathrm{V})^*\) 的连续同构。

对于所有 \(\lambda \in \mathrm{V}^*\)，族 \(\lambda^n / n!\) 在 \(\mathbf{S}(\mathrm{V})^*\) 中可求和。其和称为 \(\lambda\) 的指数，记为 \(\exp(\lambda)\)（遵循第 II 章，§6，第 1 节的约定）。设 \(x_1, \ldots, x_n \in \mathrm{V}\)；我们有

\[
\langle \exp \lambda , x _ {1} \dots x _ {n} \rangle = \frac {1}{n !} \langle \lambda^ {n}, x _ {1} \dots x _ {n} \rangle = \langle \lambda , x _ {1} \rangle \dots \langle \lambda , x _ {n} \rangle
\]

根据《代数》，第 III 章，§11，第 5 节，公式 (29)。立即可得，\(\exp(\lambda)\) 是从代数 \(\mathbf{S}(\mathrm{V})\) 到 \(k\) 的唯一同态，且延拓 \(\lambda\)。

有 \(\exp (\lambda +\mu) = \exp (\lambda)\exp (\mu)\)，对所有 \(\lambda ,\mu \in \mathrm{V}^{*}\) 都成立（第 II 章，§6，第 1 节，备注）。因此，映射 \(\exp :\mathrm{V}^{*}\to \mathbf{S}(\mathrm{V})^{*}\) 是从加法群 \(\mathrm{V}^*\) 到 \(\mathbf{S}(\mathrm{V})^{*}\) 的可逆元素乘法群的同态。族 \((\exp \lambda)_{\lambda \in \mathrm{V}^{*}}\) 是向量空间 \(\mathbf{S}(\mathrm{V})^{*}\) 中的自由族（《代数》，第 V 章，§7，第 3 节，定理 1）。

引理 1. 设 \(\Pi\) 是 \(\mathrm{V}^*\) 的一个子群，且生成向量空间 \(\mathrm{V}^*\)；\(m\) 是一个满足 \(\geq 0\) 的整数。则 \(\mathrm{pr}_m(\exp \Pi)\) 生成向量空间 \(\mathbf{S}^m (\mathrm{V}^*)\)。

根据《代数》，第 I 章，§8，第 2 节，命题 2，任意 \(m\) 个来自 \(V^{*}\) 的元素的乘积都是 \(k\)-线性组合，且组合中的元素形如 \(x^{m}\)（其中 \(x \in \Pi\)）。但 \(x^{m} = m! \operatorname{pr}_{m}(\exp x)\)。证毕。

通过结构迁移，V 的每个自同构定义 \(\mathbf{S}(\mathrm{V})\) 和 \(\mathbf{S}(\mathrm{V})^{*}\) 的自同构；这给出 \(\mathbf{GL}(\mathrm{V})\) 在 \(\mathbf{S}(\mathrm{V})\) 和 \(\mathbf{S}(\mathrm{V})^{*}\) 上的线性表示。

